\documentclass[acmsmall,nonacm,screen]{acmart}

\usepackage{booktabs}
\usepackage{array}

\citestyle{acmauthoryear}
\settopmatter{printacmref=false}

\AtEndPreamble{\theoremstyle{acmdefinition}\newtheorem{remark}[theorem]{Remark}}

\newcommand{\kw}[1]{\mathsf{#1}}
\newcommand{\tool}[2]{\texttt{#1(#2)}}
\newcommand{\any}{\textup{\texttt{\_}}}
\newcommand{\rest}{\textup{\texttt{*}}}
\newcommand{\lang}[1]{\mathcal{L}(#1)}
\newcommand{\U}{\kw{U}}
\newcommand{\fresh}{\star}
\newcommand{\meet}{\sqcap}
\newcommand{\reach}{\kw{reach}}
\newcommand{\dec}{\kw{decide}}

% A sentence that states a fact in words. The script has one check for every name.
\newcommand{\claim}[2]{#2}

% The rows of a table that the script writes.
\makeatletter
\newcommand{\rows}[1]{\@@input artifact/expected/#1.tex }
\makeatother

% written by rules.py; do not edit
\newcommand{\exActions}{5,093,247}
\newcommand{\exAlike}{5}
\newcommand{\exBlowupLargest}{eight}
\newcommand{\exBlowupNodes}{7,936}
\newcommand{\exBlowupPower}{256}
\newcommand{\exClassLists}{1,000}
\newcommand{\exClassNever}{47}
\newcommand{\exDead}{3,781}
\newcommand{\exExampleBash}{ten}
\newcommand{\exExampleRules}{thirteen}
\newcommand{\exExampleShort}{four}
\newcommand{\exFormulas}{400}
\newcommand{\exGrowthArrival}{92}
\newcommand{\exGrowthArrivalLoose}{35}
\newcommand{\exGrowthArrivalSmall}{30}
\newcommand{\exGrowthDead}{69}
\newcommand{\exGrowthKilled}{1.0}
\newcommand{\exGrowthLargest}{32}
\newcommand{\exGrowthRuns}{200}
\newcommand{\exGrowthSmallClean}{10}
\newcommand{\exLists}{2,500}
\newcommand{\exLoose}{1,459}
\newcommand{\exMasked}{2,695}
\newcommand{\exMutants}{nine}
\newcommand{\exNeeded}{6,254}
\newcommand{\exPrefix}{3,000}
\newcommand{\exRedundant}{2,283}
\newcommand{\exRules}{12,318}
\newcommand{\exSaved}{46}
\newcommand{\exSize}{800}
\newcommand{\exSpare}{6,064}
\newcommand{\exTiny}{60}
\newcommand{\exUnsat}{77}


\begin{document}

\title{A Rule Is a Clause: Dead Rules and Undecided Actions in Permission Lists}

\author{Claude-Fable-5.2}
\authornote{An AI model. The proofs in this paper are on paper and have not been checked by a
machine. A script compares every analysis with brute force over a set of actions that is
proved to be enough. At the time of writing nobody has reviewed the paper.}

\begin{abstract}
A harness decides what an agent may do with a list of rules: allow this pattern of actions,
deny that one, ask about a third. Such lists are written by hand, one rule at a time, and
they rot. A rule that other rules always beat never fires. If it was a denial, the operator
trusts a protection that does not exist, and if it was a permission, the user is asked
forever about something that the file says is allowed. An action that no rule matches falls
to a default that nobody chose. Compilers know both defects from the clauses of a match
expression, as the useless clause and the missing case, and they find both with one question:
is there a value that this pattern matches and none of those? We read a permission list as a
match expression over actions. Because the words of a command come from an open alphabet, the
question is answered in one pass over the patterns. Four diagnoses of a list follow from
polynomially many such questions (a rule is dead, a dead rule is masked, a live rule is
redundant, an action is undecided), and so does the question whether two lists decide alike.
Over a closed alphabet the first diagnosis is already coNP-complete. Compiling a list to a
decision tree gives a canonical form, which can be exponentially large. For lists of prefix
patterns it is small, and it gives the smallest list that decides the same, which is unique
up to the order of rules that do not overlap. A script compares every analysis with brute
force on \exLists{} random lists and \exActions{} actions, and it notices \exMutants{}
deliberate mistakes.
\end{abstract}

\keywords{pattern matching, usefulness, exhaustiveness, decision trees, permissions, access control, agents}

\authorsaddresses{}

\maketitle

\section{Introduction}

An agent acts through tools. It runs a command, reads a file, fetches a page. The program
around the model, the harness, decides for each such action whether it happens. The usual
means is a list of rules. A rule names a pattern of actions and a decision, and the decisions
are three: the action is allowed, it is denied, or the user is asked. One rule may say that
\tool{Bash}{git status} is allowed, another that \tool{Bash}{git push \rest} needs a
question. When several rules match an action, one of them wins, and when none does, a default
applies.

These lists are written by people, and they are written over time. A question comes up too
often, and somebody adds a permission. Something goes wrong, and somebody adds a denial.
Nobody reads the list as a whole, and nothing tells its authors what the list as a whole
decides. Two defects follow, and both are silent.

The first is a rule that never fires. Suppose that the rules are read in the order in which
they are written, that the list says early on that every \texttt{git} command needs a
question, and that somebody later adds below it that a forced push is denied. The old rule
matches every forced push first. The denial is in the file, it reads well in a review, and it
decides nothing. Harnesses that read by the kind of decision have the mirror image of this
defect. If a question beats a permission wherever the two stand, then a permission for
\texttt{git status} beside a question for every \texttt{git} command permits nothing, and the
user is asked each time. The second defect is an action that no rule matches. It gets the
default, and the default was chosen for the harness as a whole, not for this action.

Compilers have met both defects. A match expression is a list of clauses, each a pattern and
a result, and the first clause that matches is taken. A clause that no value can reach is
called useless, and a value that no clause matches shows that the match is not exhaustive.
Compilers warn about both, and since the work of \citet{Maranget2007} the warnings rest on
one question, usefulness: given some patterns and one more, is there a value that the one
matches and none of the others? A clause is useless if the answer is no for the clauses
before it. A match is not exhaustive if the answer is yes for the pattern that matches
everything.

This paper reads a permission list as a match expression over actions and asks what the
reading gives. It gives more than the two warnings, for three reasons.

\emph{The alphabet is open.} The values that a pattern of a programming language inspects are
built from the constructors of a data type, and there are finitely many. The words of a
command are not like that: any string can stand after \texttt{git}. This makes usefulness
easier, not harder. Where a pattern allows any word, a word that no rule mentions is always
the best choice, so the search never branches. One pass over the patterns answers the
question (Theorem~\ref{thm:cost}). If the alphabet is closed, as it is for a field with a
fixed set of options, the question is NP-complete, as it is for data types
(Theorem~\ref{thm:closed}).

\emph{Decisions repeat.} A match expression has a different result in every clause. A list
has three decisions, so a rule can be reached and still change nothing, because the rules
below it and the default decide the same for everything that it catches. Such a rule is
redundant. A compiler has no word for it.

\emph{Rules have an intent.} A useless clause is clutter. A dead rule may be a hole. The rule
that denies a forced push wanted its actions denied, and they are asked about. We call a dead
rule \emph{masked} if some action that it matches gets another decision than the rule's, and
\emph{loosened} if it gets a less strict one. A loosened rule is what an operator needs to
hear about first.

The paper makes these contributions.
\begin{itemize}
\item A reading of permission lists as match expressions, for both orders in which harnesses
  read rules, with a lemma that reduces every question about a list to finitely many actions
  (Section~\ref{sec:lists}).
\item Usefulness for the patterns of rules, in one pass when words are open and NP-complete
  when they are closed (Section~\ref{sec:clauses}).
\item Four diagnoses of a list and a test whether two lists decide alike, each with
  polynomially many questions of usefulness, and an account of what reading by the kind of
  decision does to the diagnoses (Section~\ref{sec:diagnoses}).
\item A decision tree for a list, whose reduced form is canonical. It can be exponentially
  large. For lists of prefix patterns it is small, and the unique smallest list is read off
  from it (Section~\ref{sec:tree}).
\item A script of about \exSize{} lines that implements all of it and compares each analysis
  with brute force (Section~\ref{sec:experiments}).
\end{itemize}

The paper does not say which actions an agent should be allowed. It also does not look below
the words of a command. A rule about the words \texttt{git push} is not a rule about pushing,
and Section~\ref{sec:limits} says what follows from that. Earlier papers of this collection
treat what an approval is \citep{tpp4}, where in a program to ask for one \citep{tpp17}, and
how a policy that remembers what happened before is learned from the answers of a user
\citep{tpp16}. This one treats the plain list, which is what harnesses have, once it is
written.

\section{A List That Rotted}\label{sec:example}

Table~\ref{tab:example} shows a list of \exExampleRules{} rules. It is invented, but each of
its defects is of a kind that the history of a real settings file produces. An action that no
rule matches is asked about. The two columns on the right say what the analyses of this paper
find, once for each of the two orders in which harnesses read such a list.

\begin{table}
\caption{A permission list and what the analyses of this paper say about each rule, if the
rules are read in the order written and the first match decides, and if a denial beats a
question and a question beats a permission wherever they stand. A pattern names a tool and
the words of its argument; \texttt{*} stands for any words that follow and \texttt{\_} for
any one word.}\label{tab:example}
\begin{tabular}{rllll}
\toprule
 & pattern & decision & in the order written & by the kind of decision \\
\midrule
\rows{example}
\bottomrule
\end{tabular}
\end{table}

\paragraph{In the order written.} Rule~2 asks about every \texttt{git} command. Rule~3 was
added later, after a forced push destroyed somebody's work, and it was added below rule~2. It
is dead: rule~2 matches every forced push before rule~3 is looked at. The author wanted a
denial and got a question, which an agent that runs with its questions answered in advance
will pass. Rule~11 is the same mistake with a worse outcome. Rule~10 allows reading
everything under \texttt{src}, the denial for the secrets below it never fires, and the
secrets are read without a question. Rule~6 and the last rule are dead as well, but
harmlessly: rule~6 repeats rule~4, and the last rule asks where rule~2 already asks. Rules~5
and~8 are alive and change nothing. They ask, and without them the default would ask.

\claim{example}{The analyses report exactly this. Rules 3, 6 and 11 and the last rule are
dead. Rules 3 and 11 are masked, and both are loosened. Rules 5 and 8 are redundant, and both
tools have undecided actions. A forced push is asked about, and a secret is read. The
\exExampleBash{} rules for \texttt{Bash} decide like a list of \exExampleShort{}.} The four
are \tool{Bash}{git status}, \tool{Bash}{npm test} and \tool{Bash}{ls \rest}, which allow,
and \tool{Bash}{rm -rf \rest}, which denies. Everything else that the ten rules say is what
the default says. A reader of the four sees at once that no rule mentions a push. A reader of
the ten does not.

\paragraph{By the kind of decision.} Now let every denial beat every question, and every
question beat every permission. The two denials fire, and the holes are closed. Other rules
die instead. Rule~1 allows \texttt{git status}, and the question of rule~2 beats it. Rules~4
and~6 allow \texttt{npm test}, and the question of rule~5 beats them. The user is asked about
both commands every time, and the file says that they are allowed. Rule~5, which changed
nothing before, now decides that \texttt{npm test} is asked about.

\claim{example-classes}{Read in this way, rules 1, 4 and 6 and the last rule are dead, and
the first three of them are masked. No rule is loosened. Rule~8 alone is redundant. A forced
push and the reading of a secret are denied, and \texttt{git status} and \texttt{npm test}
are asked about.}

Neither order is free of dead rules. The first kills denials, which is dangerous. The second
kills permissions, which is costly, because a user who is asked about what is allowed learns
to answer without reading.

\section{Lists}\label{sec:lists}

We fix a set of \emph{tools} and a set $\mathcal{W}$ of \emph{words}. The set of words is
infinite: it holds every string that can stand in a command or a path. An \emph{action}
$t(w_1 \ldots w_n)$ is a tool $t$ and a sequence of $n \ge 0$ words. For a shell command the
words are the words of the command, and for a file they are the parts of its path.

A \emph{pattern} $t(\pi_1 \ldots \pi_k)$ or $t(\pi_1 \ldots \pi_k\,\rest)$ is a tool and a
sequence of $k \ge 0$ \emph{items}, each a word or the wildcard~\any, and it may end
in~\rest. The pattern $t(\pi_1 \ldots \pi_k)$ matches the action $t(w_1 \ldots w_n)$ if $n =
k$ and every item $\pi_i$ is $w_i$ or \any. The pattern $t(\pi_1 \ldots \pi_k\,\rest)$
matches it if $n \ge k$ and the same holds for the first $k$ words. We write $\lang{p}$ for
the set of actions that $p$ matches and $|p|$ for its number of items, and we say that $p$
\emph{has a rest} if it ends in~\rest. Every pattern matches some action. A \emph{prefix
pattern} is a pattern without a wildcard.

A \emph{rule} is a pattern and a \emph{decision}, one of $\kw{allow}$, $\kw{ask}$ and
$\kw{deny}$. We order the decisions by strictness, $\kw{allow} < \kw{ask} < \kw{deny}$. A
\emph{list} $L$ is a sequence of rules $(p_1, d_1), \ldots, (p_n, d_n)$ and a default
decision $d_0$. It decides an action $a$ by the first rule that matches:
\[
  \dec(L, a) \;=\; d_i \quad\text{for the least $i$ with } a \in \lang{p_i},
  \qquad\text{and } d_0 \text{ if there is none.}
\]
The actions that reach rule $i$ are $\reach(i) = \lang{p_i} \setminus \bigcup_{j < i}
\lang{p_j}$.

\paragraph{Two orders.} Harnesses differ in which of several matching rules wins. Some read
the rules in the order in which they are written, as a match expression and a firewall do.
Others go by the kind of decision: a denial beats a question, and a question beats a
permission, wherever the rules stand and however specific they are. Claude Code documents
the second order \citep{claudecodepermissions}. It is the first order applied to another
list, the one that has been sorted by decision, the strictest first, with a stable sort. We
call that list the list \emph{in the order of classes}. Everything below is said about lists
read by the first match and so applies to both, and Section~\ref{sec:diagnoses} says what the
sorting changes.

\paragraph{Finitely many actions are enough.} A list is a statement about infinitely many
actions. The following lemma makes it a statement about finitely many, which is what lets a
script check the rest of the paper by brute force.

\begin{lemma}[A finite test set]\label{lem:finite}
Let $P$ be a finite set of patterns, $W$ the set of words that occur in them, $m$ the
greatest number of items of a pattern of $P$, and $\fresh$ a word that is not in $W$. For an
action $a = t(w_1 \ldots w_n)$ let $\bar{a} = t(w'_1 \ldots w'_{n'})$ with $n' = \min(n, m +
1)$, and $w'_i = w_i$ if $w_i \in W$ and $w'_i = \fresh$ otherwise. Then every pattern of $P$
matches $a$ exactly if it matches $\bar{a}$.
\end{lemma}

\begin{proof}
Let $p \in P$ have $k \le m$ items. An item that is a word is in $W$, so it equals $w_i$
exactly if it equals $w'_i$, and a wildcard matches both. It remains to compare the lengths.
If $p$ has a rest, it needs $n \ge k$, and $\min(n, m + 1) \ge k$ says the same because $k <
m + 1$. If it has none, it needs $n = k$, and $\min(n, m + 1) = k$ says the same for the same
reason.
\end{proof}

So two lists that agree on all actions with at most $m + 1$ words over $W \cup \{\fresh\}$
agree on all actions of the tools that they mention, where $W$ and $m$ are taken over the
patterns of both. On every other tool they agree exactly if their defaults are equal.

\section{Clauses}\label{sec:clauses}

Read a list as a match expression. The value that is matched is the action, a rule is a
clause, its pattern is the pattern of the clause, and its decision is the result. The tool is
a constructor at the head, and the words form a list, which a pattern takes apart from the
left: a word item is a literal, \any{} is the wildcard for one element, an end without a rest
is the pattern for the empty list, and \rest{} is the wildcard for the remaining list.

\subsection{Usefulness}

Let $P$ be a sequence of patterns, the \emph{rows}, and $q$ a pattern. We say that $q$ is
\emph{useful} against $P$ if some action is matched by $q$ and by no row, and we call such an
action a \emph{witness}. Rows for another tool than that of $q$ match no action that $q$
matches, so we drop them and the tool, and we write a row as its items $r_1 \ldots r_k$.

The function $\U_i(R, q)$ decides usefulness after $i$ words have been chosen. Its argument
$R$ is the set of rows that can still match: the rows with at least $i$ items whose first $i$
items match the chosen words. It returns a witness or $\bot$, and it is defined by cases, of
which the first that applies is taken.
\begin{enumerate}
\item If some row of $R$ has exactly $i$ items and a rest, the result is $\bot$.
\item If $i \ge |q|$ and no row of $R$ has exactly $i$ items, the result is the action of
  the $i$ chosen words.
\item If $i \ge |q|$ and $q$ has no rest, the result is $\bot$.
\item If $i \ge |q|$, or if $q_{i+1}$ is the wildcard, the next word is $\fresh$, a word that
  occurs in no pattern, and the result is $\U_{i+1}(\{r \in R : |r| > i \text{ and } r_{i+1}
  = \any\},\, q)$.
\item Otherwise $q_{i+1}$ is a word $s$, the next word is $s$, and the result is
  $\U_{i+1}(\{r \in R : |r| > i \text{ and } r_{i+1} \in \{s, \any\}\},\, q)$.
\end{enumerate}
We write $\U(P, q)$ for $\U_0(R, q)$, where $R$ is the set of rows of $P$ for the tool of
$q$.

In case~1 a row has matched all its items and takes whatever follows, so nothing is left for
$q$. In case~2 the action may end here, and no row wants an action that ends here. In case~3
the action must end here, because $q$ says so, and some row wants it. Cases~4 and~5 read one
more word.

\begin{theorem}[Usefulness]\label{thm:useful}
$\U(P, q)$ returns a witness if $q$ is useful against $P$, and $\bot$ if it is not.
\end{theorem}

\begin{proof}
The computation ends, because a row leaves $R$ at the latest when $i$ exceeds its number of
items, and with no rows left and $i \ge |q|$ case~2 applies.

Fix $i$ chosen words, and let $A$ be the set of actions that begin with them and that $q$
matches. We show the following by induction on the number of steps that are left. Suppose
that $A$ is not empty and that $R$ is the set of all rows that match some action beginning
with the chosen words, each of them with at least $i$ items. Then $\U_i(R, q)$ returns an
element of $A$ that no row matches if there is one, and $\bot$ otherwise. At the start this
is the theorem: no words are chosen, $A$ is $\lang{q}$, and every row matches some action.

In case~1 the row matches every action that begins with the chosen words. In case~2 the
action of the chosen words is in $A$: the set $A$ is not empty, so $q$ matches the chosen
words on its first $|q|$ items, and it has a rest if $i > |q|$. A row that matches this
action is in $R$ and has exactly $i$ items, and there is none. In case~3 that action is the
only element of $A$, and a row of $R$ with exactly $i$ items matches it.

In cases~4 and~5 the action of the chosen words is not in $A$, or it is matched by a row as
in case~3. So a witness has at least one more word. A row that matches an action with more
than $i$ words has more than $i$ items, because the rows of $R$ with exactly $i$ items have
no rest. In case~5 the next word of every action of $A$ is $s$, and the rows that match some
action beginning with the longer sequence are those named. In case~4 the next word can be any
word. Suppose that a witness $a$ has the word $w$ there, and let $a'$ be $a$ with $\fresh$ in
place of $w$. Then $q$ matches $a'$, because $q$ has a wildcard or a rest at this place. A
row that matches $a'$ has a wildcard at this place, because $\fresh$ occurs in no pattern, so
it matches $a$ too, and there is no such row. So $a'$ is a witness as well, and it is enough
to look for witnesses that continue with $\fresh$. The rows that match some action beginning
with the chosen words and $\fresh$ are those named. In both cases the longer sequence
satisfies the suppositions, and the induction hypothesis gives the claim.
\end{proof}

The step from $a$ to $a'$ is the whole difference to data types. For a type with finitely
many constructors there is no fresh one, and the search must try them all.
\citet{Maranget2007} makes the same step for a signature that the patterns do not exhaust.
Here that is always so.

\begin{theorem}[One pass]\label{thm:cost}
$\U(P, q)$ takes time linear in the size of $P$ and $q$, counted in items.
\end{theorem}

\begin{proof}
There is one step for every position, up to one more than the longest row and $|q|$. In a
step each row of $R$ is looked at once, and a row with $k$ items is in $R$ for at most $k +
1$ steps.
\end{proof}

\subsection{The Meet of Two Patterns}

The diagnoses of Section~\ref{sec:diagnoses} ask about actions that two rules both match.
These actions are again the actions of a pattern. Let $p$ and $q$ be patterns for the same
tool with $|p| \le |q|$. Their \emph{meet} $p \meet q$ is undefined if at some position both
have a word and the words differ, or if $|p| < |q|$ and $p$ has no rest. Otherwise it has the
items of $q$, with the word of $p$ wherever $q$ has a wildcard and $p$ has a word, and it has
a rest if both have one. Patterns for different tools have no meet.

\begin{lemma}[Meet]\label{lem:meet}
If $p \meet q$ is defined, then $\lang{p \meet q} = \lang{p} \cap \lang{q}$. If it is not,
the intersection is empty.
\end{lemma}

\begin{proof}
An action that both match has the word of either pattern wherever one of them has a word, so
two different words exclude it. Its length is $|p|$ if $p$ has no rest, which $q$ accepts
only if $|q| = |p|$. Conversely, an action with the words of both and a length that both
accept is matched by both, and these are the conditions that $p \meet q$ imposes.
\end{proof}

\subsection{A Closed Alphabet}

Not every field of an action is open. A tool may take options from a fixed set, such as a
sandbox that is on or off. Suppose that the words come from a finite set that is given with
the list. Then Lemma~\ref{lem:finite} holds without the fresh word, and usefulness is still
decidable, by trying every word in case~4. The single pass is lost for good.

\begin{theorem}[Closed alphabets]\label{thm:closed}
If the words come from a given set of two words, it is coNP-complete to decide whether a rule
of a list is dead, that is, whether no action reaches it.
\end{theorem}

\begin{proof}
A rule that is not dead has a witness with at most one word more than the longest pattern, by
the argument of Lemma~\ref{lem:finite}, and a witness is checked in polynomial time. So it is
in NP to decide that a rule is not dead.

For the hardness, take a formula in conjunctive normal form with variables $x_1, \ldots,
x_n$, in which no clause contains a variable and its negation, and let the words be
$\texttt{t}$ and $\texttt{f}$. An assignment is an action of $n$ words. For every clause of
the formula the list has a rule whose pattern has $n$ items and no rest: the item at position
$j$ is $\texttt{f}$ if the clause contains $x_j$, it is $\texttt{t}$ if the clause contains
the negation of $x_j$, and it is the wildcard otherwise. This pattern matches exactly the
assignments that make the clause false. The last rule of the list has $n$ wildcards and no
rest. An action reaches it exactly if it is an assignment that makes no clause false. So the
last rule is dead exactly if the formula cannot be satisfied, and satisfiability is
NP-complete \citep{Cook1971}.
\end{proof}

With an open alphabet the same list is harmless. The action of $n$ fresh words reaches the
last rule, whatever the formula is.

\section{Four Diagnoses}\label{sec:diagnoses}

Let $L$ be a list with rules $(p_1, d_1), \ldots, (p_n, d_n)$ and default $d_0$. We write
$P_{<i}$ for the patterns of the rules before rule $i$, and $P \setminus i$ for a sequence of
patterns without that of rule $i$.

\begin{itemize}
\item Rule $i$ is \emph{dead} if $\reach(i)$ is empty, and \emph{alive} otherwise.
\item A dead rule $i$ is \emph{masked} if $\dec(L, a) \ne d_i$ for some $a \in \lang{p_i}$,
  and \emph{loosened} if $\dec(L, a) < d_i$ for some such $a$.
\item A live rule $i$ is \emph{redundant} if the list without it decides every action as $L$
  does.
\item An action is \emph{undecided} if no rule matches it.
\end{itemize}

A dead rule that is not masked repeats what the rules above it say. A masked rule contradicts
them and loses. A loosened rule loses to a weaker decision: it asked for a denial where the
list asks or allows, or for a question where the list allows. The opposite mask, a permission
that never permits, costs the user a question or the agent an action. A loosened rule costs
the protection that its author believes in.

\begin{theorem}[The diagnoses]\label{thm:diagnoses}
For a list with $n$ rules:
\begin{enumerate}
\item Rule $i$ is dead exactly if $\U(P_{<i}, p_i) = \bot$.
\item For $j < i$, some action that $p_i$ matches is decided by rule $j$ exactly if $p_i
  \meet p_j$ is defined and useful against $P_{<j}$. So a dead rule $i$ is masked exactly if
  this holds for some $j < i$ with $d_j \ne d_i$, and loosened exactly if it holds for some
  $j < i$ with $d_j < d_i$.
\item A live rule $i$ is redundant exactly if neither of the following holds: $d_0 \ne d_i$
  and $p_i$ is useful against $P \setminus i$; or there is a $k > i$ with $d_k \ne d_i$ such
  that $p_i \meet p_k$ is defined and useful against $P_{<k} \setminus i$.
\item A tool $t$ has an undecided action exactly if $t(\rest)$ is useful against the patterns
  of $L$.
\end{enumerate}
All four are decided with $O(n^2)$ questions of usefulness.
\end{theorem}

\begin{proof}
Parts~1 and~4 restate the definitions. For part~2, the actions that $p_i$ and $p_j$ both
match are those of $p_i \meet p_j$ by Lemma~\ref{lem:meet}, and rule $j$ decides such an
action exactly if no rule before $j$ matches it. An action of a dead rule is decided by an
earlier rule, the first that matches it, which gives the rest of part~2.

For part~3, removing rule $i$ changes the decision only of actions in $\reach(i)$. Such an
action is then decided by the first later rule that matches it, or by the default. It goes to
the default exactly if no other rule matches it, and an action with this property exists
exactly if $p_i$ is useful against $P \setminus i$. It goes to rule $k > i$ exactly if it is
in $\lang{p_i \meet p_k}$ and no rule before $k$ other than $i$ matches it. The rule is
redundant exactly if no action of $\reach(i)$ gets another decision than $d_i$ in this way.
\end{proof}

Part~4 says whether a tool has undecided actions and gives one. Section~\ref{sec:tree}
describes all of them.

\begin{remark}[Exceptions]
A live rule can lose some of its actions to an earlier rule with another decision, by the
question of part~2. We do not call that a defect. It is how an exception is written: a
permission for \texttt{git status} above a question for every \texttt{git} command takes one
command out of the question, on purpose. A dead rule is an exception to nothing. Part~2 can
still be asked for a live denial, and its witnesses then list what the denial does not deny.
\end{remark}

\paragraph{What the order of classes does.} Sorting a list by decision is a policy about
masks.

\begin{proposition}[Classes]\label{prop:classes}
In a list in the order of classes no rule is loosened.
\end{proposition}

\begin{proof}
A rule before rule $i$ has a decision that is at least as strict as $d_i$. An action that
$p_i$ matches and that does not reach rule $i$ is decided by such a rule.
\end{proof}

So the order of classes removes the dangerous defect by construction, and the forced push of
Section~\ref{sec:example} is denied under it. It pays in two ways. Rules still die, and the
ones that die are permissions and questions: a permission inside a region that is asked about
or denied permits nothing, and its author is not told. And an exception to a denial cannot be
written at all, because the denial of the larger region wins wherever the exception stands.
The diagnoses apply to the sorted list as they are, and Section~\ref{sec:experiments} counts
how many permissions of a random list survive the sorting.

\paragraph{Two lists.} The same questions decide whether a change to a list changes what it
decides.

\begin{theorem}[Equivalence]\label{thm:equivalence}
Let $L$ and $L'$ be lists with rules $(p_i, d_i)$ and $(p'_j, d'_j)$, defaults $d_0$ and
$d'_0$, and patterns $P$ and $P'$. Some action is decided differently by the two exactly if
one of the following is useful:
\begin{enumerate}
\item $p_i \meet p'_j$ against $P_{<i}$ and $P'_{<j}$ together, for some $i$ and $j$ with
  $d_i \ne d'_j$;
\item $p_i$ against $P_{<i}$ and all of $P'$, for some $i$ with $d_i \ne d'_0$, or the same
  with the two lists exchanged;
\item $t(\rest)$ against $P$ and $P'$ together, for some tool $t$, if $d_0 \ne d'_0$.
\end{enumerate}
Every witness is an action on which the lists differ. So whether two lists with $n$ and $n'$
rules decide alike is decided with $O(n\,n')$ questions of usefulness.
\end{theorem}

\begin{proof}
An action is decided in each list by its first matching rule or by the default. The three
cases are the three ways in which the two outcomes can be a rule and a rule, a rule and a
default, and two defaults, and in each the question describes the actions with this pair of
outcomes, by Lemma~\ref{lem:meet} and the definition of usefulness.
\end{proof}

Part~3 of Theorem~\ref{thm:diagnoses} is the case in which $L'$ is $L$ without one rule.

\section{One Tree for Every Question}\label{sec:tree}

The questions of Section~\ref{sec:diagnoses} are asked rule by rule. A compiler does not
answer a match expression clause by clause. It compiles the clauses to a decision tree that
inspects every part of the value once \citep{Augustsson1985,Maranget2008}. The tree of a
permission list is an object of its own interest, because it shows what the list decides
without showing how the list was written.

\subsection{The Tree of a List}

Fix a tool. A node of the tree stands for a sequence of words that has been read. It is built
from a set of \emph{rows}, the rules for the tool that match these words on their first
items, and from a \emph{floor}, the first rule that has matched all its items and has a rest.
A rule after the floor decides nothing below the node, so the rows are only the rules before
the floor. The node has three parts:
\begin{itemize}
\item the rule that decides if the action \emph{ends here}: the first among the floor and the
  rows with exactly as many items as words have been read, or no rule;
\item a child for every word that some row has as its next item, built from the rows whose
  next item is that word or the wildcard;
\item a child for \emph{every other word}, built from the rows whose next item is the
  wildcard.
\end{itemize}
A node without rows that have a next item is a \emph{leaf} if the floor also decides for the
end. A leaf names the floor, or no rule, for everything that begins with the words read. The
root is built from all rules for the tool, with no words read and no floor. The tree of a
tool that the list does not mention is the leaf of no rule.

\begin{theorem}[The tree]\label{thm:tree}
Following an action through the tree of its tool, by the child for its next word if there is
one and by the child for other words if not, ends at the rule that decides the action, or at
no rule if the action is undecided. The rules that occur in the trees of a list are its live
rules.
\end{theorem}

\begin{proof}
By induction along the walk, the rows of a node are the rules before the floor that match the
words read so far on their first items, and the floor is the first rule that matches every
action beginning with these words by having a rest. The first rule that matches the action is
the first among the floor and the rows that match all of it. When the action ends, these rows
are the ones with exactly as many items as the action has words, because a row with fewer
items and a rest would be the floor or come after it. So the walk ends at the first matching
rule. Every rule in a tree is reached by the action that leads to it, read with fresh words
on the edges for other words, and conversely the walk of a witness of a live rule ends at
that rule.
\end{proof}

For a list of prefix patterns the tree has at most one node for every prefix of a pattern,
and one leaf below each, so it is linear in the size of the list.

\subsection{A Canonical Form}

Two trees can differ and decide alike: they name different rules for the same decision, and
they can have a child for a word under which the same happens as under every other word. The
\emph{reduced tree} removes both. It has decisions in place of rules, with the default where
no rule decides. It is built from the leaves up. A child for a word is dropped if its reduced
tree equals that of the child for other words. A node that has no child for a word left, and
whose child for other words is a leaf with the node's own decision for the end, is replaced
by that leaf.

\begin{theorem}[Canonical form]\label{thm:canonical}
Two lists decide every action alike exactly if, for every tool, their reduced trees are
equal.
\end{theorem}

\begin{proof}
Equal trees decide alike by Theorem~\ref{thm:tree}, since the reduction does not change the
decision at which a walk ends. For the converse, call a tree with decisions \emph{reduced} if
no child for a word equals the child for other words beside it, and no node could be replaced
by a leaf as above. We show that two reduced trees that represent the same function $f$ from
sequences of words to decisions are equal, by induction on the size of the two together.

For a word $z$ let $f_z$ be the function that maps $u$ to $f(z\,u)$. In a node, all words
without a child have the same $f_z$, and these are all words but finitely many. Because the
alphabet is infinite, at most one function $g$ equals $f_z$ for all but finitely many $z$, so
$g$ is determined by $f$, and the child for other words represents it.

Suppose that one tree is a leaf, so that $f$ is constant, and the other is a node. The
children of the node represent the same constant function, so by induction all of them equal
the leaf. Then the node has no child for a word, since such a child would equal the child for
other words, and it could be replaced by the leaf, which contradicts that the tree is
reduced. So if one tree is a leaf, both are, and they are equal. Now let both be nodes. Their
decisions for the end are both $f$ of the empty sequence. Their children for other words both
represent $g$ and are equal by induction. A word $z$ without a child has $f_z = g$. A word
with a child has $f_z \ne g$, because the child differs from the child for other words, and by
induction two reduced trees that differ represent different functions. So both nodes have
children for the same words, the words with $f_z \ne g$, and the two children for a word
represent the same function and are equal by induction.
\end{proof}

The proof uses the open alphabet once more. Over a closed alphabet there is no child for
other words to compare with, and a canonical form needs a fixed order of the words, as the
decision diagrams of \citet{Bryant1986} need a fixed order of the variables.

The reduced tree answers what Theorem~\ref{thm:diagnoses} left open. If ``no rule'' is kept
as a fourth decision, the paths that end there describe all undecided actions of a tool, each
as a sequence of words, with the words that are excluded where the path takes a child for
other words. A change to a list changes what it decides exactly if it changes a reduced tree,
and the subtrees that differ are the actions whose decision changed. This is the
change-impact analysis of \citet{Fisler2005} for this kind of policy.

\paragraph{The price.} A canonical form is not free. Take $n$ rules, each with $2n$ items
and no rest, where rule $j$ has the word \texttt{a} at the places $j$ and $n + j$ and
wildcards elsewhere. All of them allow, and the default denies. After $n$ words the tree must
know at which of the first $n$ places an \texttt{a} stood, because every such set of places
allows other continuations. So among the first $n$ levels no child for \texttt{a} equals the
child beside it, and the reduced tree has at least $2^n$ nodes. \claim{blowup}{The script
builds these trees for $n$ up to \exBlowupLargest{}. The largest has \exBlowupNodes{} nodes,
where $2^n$ is \exBlowupPower{}, and each tree has more than twice the nodes of the one
before. The questions of Theorems~\ref{thm:diagnoses} and~\ref{thm:equivalence} answer for
the same lists without a tree: no rule is dead or redundant, the list read backwards decides
alike, and the list without its first rule does not.} Decision trees for the patterns of
programming languages have the same property \citep{Sekar1995}. So the tree is the right
means for showing a list and for comparing many lists with one, and the questions of
Section~\ref{sec:diagnoses} are the right means for a single answer.

\subsection{The Smallest List of Prefix Patterns}

For prefix patterns the reduced tree is small, and it is close to a list already. No rule has
a wildcard, so the child for other words is built from no rows and is a leaf. Call its
decision the decision \emph{below} the node, and for a leaf call its own decision the
decision below it. The \emph{normal form} of a list of prefix patterns has, for every node of
the reduced tree with prefix $w$ and tool $t$,
\begin{itemize}
\item the rule $t(w\,\rest)$ with the decision below the node, if that differs from the
  decision below the parent, or from the default at the root; and
\item the rule $t(w)$ with the node's decision for the end, if that differs from the decision
  below the node.
\end{itemize}
A rule for a longer prefix stands before a rule for a shorter one, and $t(w)$ before
$t(w\,\rest)$.

\begin{theorem}[Normal form]\label{thm:normal}
Let $L$ be a list of prefix patterns and $N$ its normal form.
\begin{enumerate}
\item $N$ decides like $L$.
\item Every list of prefix patterns that decides like $L$ contains every pattern of $N$.
\item So no list of prefix patterns that decides like $L$ has fewer rules than $N$. One with
  as many has the rules of $N$, and its order differs from that of $N$ only between rules
  that match no common action.
\end{enumerate}
\end{theorem}

\begin{proof}
\emph{Part 1.} Call a node \emph{marked} if $N$ has a rule $t(w\,\rest)$ for it. Going up
from a node to the nearest marked node at or above it, the decision below does not change, by
the definition of a mark, and if there is no marked node it is the default. Let $t(u)$ be an
action, and let $w$ be the longest prefix of $u$ that has a node. The rules of $N$ that match
the action are $t(u)$ itself, if it is a rule, and the rules $t(v\,\rest)$ of the marked
nodes at or above the node of $w$, the nearest of them first. If $u$ is longer than $w$, or
the node of $w$ is a leaf, the tree says the decision below that node, and so does the first
matching rule, or the default if there is none. If $u$ is $w$, the tree says the node's
decision for the end. If that differs from the decision below, $t(w)$ is a rule, stands
first, and says it. If not, the case is as before.

\emph{Part 2.} Let $L'$ be a list of prefix patterns that decides like $L$, and let $z$ be a
word that occurs in neither list. Because $z$ occurs in no pattern, the patterns of $L'$ that
match an action $t(u\,z)$ are those of the form $t(v\,\rest)$ with $v$ a prefix of $u$.
Consider a rule $t(w\,\rest)$ of $N$. If $w$ is empty, the action $t(z)$ gets the decision
below the root, which is not the default, and only $t(\rest)$ can match it, so $t(\rest)$ is
in $L'$. Otherwise $w$ is $v\,s$ for a word $s$. The action $t(w\,z)$ gets the decision below
the node of $w$, and $t(v\,z)$ gets the decision below its parent, and these differ. If
$t(w\,\rest)$ were not in $L'$, the same patterns of $L'$ would match both actions, the first
of them would decide both, and the two would get the same decision. Now consider a rule
$t(w)$ of $N$. The actions $t(w)$ and $t(w\,z)$ get different decisions. The patterns of $L'$
that match $t(w\,z)$ match $t(w)$ as well, and the only other pattern that can match $t(w)$
is $t(w)$ itself. If it were not in $L'$, both actions would get the same decision.

\emph{Part 3.} The patterns of $N$ are different from one another, so by part~2 a list $L'$
that decides like $L$ has at least as many rules, and one with as many has exactly the
patterns of $N$. No rule of such a list is dead, since the list without a dead rule would
decide alike with fewer rules. So $t(w)$ stands before every $t(v\,\rest)$ with $v$ a prefix
of $w$, and $t(w\,\rest)$ before every $t(v\,\rest)$ with $v$ a shorter prefix of $w$,
because in the other order the first would be dead. These are all the pairs of patterns that
match a common action, and $N$ has them in the same order. Then $t(w)$ decides the action
$t(w)$ and $t(w\,\rest)$ decides $t(w\,z)$ in both lists, so the decisions of the rules are
those of $N$.
\end{proof}

A smallest list is thus not found by search. It is forced, and it is forced by the open
alphabet: the word $z$ that no rule mentions pins every pattern down. For routing tables,
whose alphabet has two letters, the decision that a prefix carries is a choice, and the
smallest table needs the algorithm of \citet{Draves1999}, which carries sets of candidates up
the tree. The normal form is also what a list should look like after a cleanup. It has no
dead and no redundant rule, and two lists of prefix patterns decide alike exactly if their
normal forms are the same list up to the order of rules that do not overlap. For lists with
wildcards we know no such form, and we do not know how hard the smallest list is to find.

\section{Experiments}\label{sec:experiments}

The script \texttt{rules.py} implements Sections~\ref{sec:lists} to~\ref{sec:tree}. It trusts
one function, the one that says whether a pattern matches an action, and compares everything
else with it by brute force over the actions of Lemma~\ref{lem:finite}. It needs nothing but
Python, and \texttt{make -C artifact check} runs it and compares its output with the numbers
in this paper.

\paragraph{Against brute force.} The script draws \exLists{} random lists of two to eight
rules, \exRules{} rules in all. A pattern follows a path through a small tree of commands
(\texttt{git push origin main}, \texttt{npm test}, \texttt{rm -rf} and the like), stops after
one to four words, has a wildcard in place of a word with some probability, and has a rest
more often than not. \claim{force}{For every list the script tries every action of the finite
test set, \exActions{} in all, and finds that usefulness with its witness, the meet, the four
diagnoses with the loosened rules, the tree, and both tests of equivalence agree with what
the definitions say. Of the rules, \exDead{} are dead, \exMasked{} of them masked and
\exLoose{} of those loosened, and \exRedundant{} are redundant. Each list is compared with
every list that lacks one of its rules: the two decide alike for \exSpare{} rules, the dead
and the redundant ones, and differently for \exNeeded{}. Each list is also compared with a
second random list, and they decide alike in \exAlike{} cases.} The high share of dead rules
says something about the generator, which draws from few commands, and nothing about real
lists.

\paragraph{The normal form.} \claim{normal}{For \exPrefix{} random lists of prefix patterns,
the normal form decides like the list, is its own normal form, has no pattern that the list
has not, and has \exSaved{} percent fewer rules in total. For \exTiny{} lists with at most
three rules in normal form, a search over lists of prefix patterns with fewer rules finds
none that decides alike.} The search is over every sequence of rules whose patterns use the
words and lengths of the list, and it is feasible only for the smallest cases.

\paragraph{Lists that grow.} How fast does a list rot? The script writes lists of four to
\exGrowthLargest{} rules, each rule drawn as above without wildcards, \exGrowthRuns{} lists
of each length, with a default that asks. Because the rules are drawn independently, it makes
no difference to the list as a whole whether a new rule is put last, as by an author who
appends to a file, or first, as by an author who wants the new rule to win. It makes a
difference to the new rule. Put last, it can be dead on arrival. Put first, it is alive, and
it can kill rules that were alive before. Table~\ref{tab:growth} shows both.

\begin{table}
\caption{Lists written one random rule at a time. Left: the percentage of rules with each
diagnosis, and the percentage of lists with no dead and no redundant rule. Right: the
percentage of lists in which the rule that was put last is dead, or dead and loosened, and
the mean number of live rules that the rule that was put first kills.}\label{tab:growth}
\begin{tabular}{r@{\qquad}rrrrr@{\qquad}rr@{\qquad}r}
\toprule
 & \multicolumn{4}{c}{rules} & lists & \multicolumn{2}{c}{the last rule} & the first \\
\cmidrule(r){2-5}\cmidrule(lr){7-8}
rules & dead & masked & loosened & redundant & clean & dead & loosened & kills \\
\midrule
\rows{growth}
\bottomrule
\end{tabular}
\end{table}

\claim{growth}{The shares of dead, of masked and of loosened rules grow with the length of
the list, and more than half of the dead rules are masked at every length. At
\exGrowthLargest{} rules \exGrowthDead{} percent of the rules are dead and no list is clean.
A rule that is put last is dead on arrival in \exGrowthArrivalSmall{} percent of the lists of
four rules and in \exGrowthArrival{} percent of the lists of \exGrowthLargest{}, and in
\exGrowthArrivalLoose{} percent of the lists of that length it is loosened as well. A rule that is put first kills more
rules the longer the list is, \exGrowthKilled{} on average at that length.} Of the lists of
four rules, \exGrowthSmallClean{} percent are clean. These numbers come from random rules
over a dozen commands and are far above what a careful author produces. The direction is
what carries over. A rule that is added without looking at the list is more likely to be dead
the longer the list is, and the author who appends loses the rule that was just meant to
change something.

\paragraph{The order of classes.} \claim{classes}{The script sorts \exClassLists{} random
lists of eight rules by decision. No rule of a sorted list is loosened, as
Proposition~\ref{prop:classes} says, and \exClassNever{} percent of the permissions are
dead.} Again the share is a property of the generator. The point is that the order of classes
moves the defect from denials to permissions and does not remove it.

\paragraph{A closed alphabet.} \claim{closed}{For \exFormulas{} random formulas with three
to seven variables, of which \exUnsat{} cannot be satisfied, the last rule of the list of
Theorem~\ref{thm:closed} is dead exactly for those when the words are \texttt{t} and
\texttt{f} only, and it is never dead when the alphabet is open.}

\paragraph{Planted mistakes.} A check that never fails checks nothing. The script has
\exMutants{} switches, each of which plants one mistake in an analysis: usefulness that lets
a row that ends take what follows, that forgets a row with a rest, or that keeps the rows
that ask for a known word where a fresh one is chosen; a meet that ignores lengths;
redundancy that forgets the default; a test of equivalence that forgets the rules of the
second list that the first leaves to its default; a tree whose floor is the last rule instead
of the first; a reduction that keeps every child; a normal form without rules for a prefix
alone. \claim{all-noticed}{Each of them is noticed by a comparison that fails.}

\section{Related Work}\label{sec:related}

\paragraph{Pattern matching.} The warnings of compilers for useless clauses and missing cases
are as old as pattern matching. \citet{Maranget2007} gives the algorithm on which this paper
rests and proves it correct, and \citet{Karachalias2015} extend the question to guards and
indexed types. Compilation to decision trees goes back to \citet{Augustsson1985}, and
\citet{Maranget2008} studies how to keep the trees small. \citet{Sekar1995} show that such
trees can be exponential and that the order of inspection matters. The reduction of
Theorem~\ref{thm:closed} is the standard one for patterns over a finite signature. What is
new here is the object, and with it the open alphabet, which turns the hard question into an
easy one and forces the smallest list, and the diagnoses that a match expression has no use
for. A clause of a program is not redundant, because clauses have different results, and it
has no intent that a mask could defeat.

\paragraph{Firewalls and access control.} Lists of rules with a first match decide which
packets pass a firewall, and their defects have been classified. \citet{AlShaer2004} define
shadowing, correlation, generalization and redundancy for pairs of rules, and
\citet{Yuan2006} check a rule against all rules before it with decision diagrams. A packet
has a fixed number of fields over finite ranges, and that is the closed case of this paper.
Our dead rule is a shadowed rule, taken against all earlier rules together, which matters
here as there: the rule $t(w\,\rest)$ is dead below $t(w)$ and $t(w\,\any\,\rest)$, and
neither of the two covers it alone. The masked and the loosened rule refine the shadowed one
by what its author wanted. \citet{Fisler2005} verify access-control policies and compute the
impact of a change with decision diagrams, as Section~\ref{sec:tree} does with reduced trees.
\citet{Draves1999} find the smallest routing table, and the comparison with
Theorem~\ref{thm:normal} shows what the open alphabet buys. That a default should deny is one
of the principles of \citet{Saltzer1975}. The undecided actions are where that principle
applies, and a harness that asks by default follows it in a weaker form.

\paragraph{Decision lists.} A list of rules with a first match and a default is a decision
list in the sense of \citet{Rivest1987}, who studies how to learn one from examples. The
diagnoses of this paper apply to a learned list as to a written one, since a learner that
appends a rule can append a dead one.

\paragraph{Approvals.} Other papers of this collection treat the approval that a rule asks
for: as a resource that an action consumes \citep{tpp4}, as something that only actions under
untrusted influence need \citep{tpp7}, and as a computation to be placed in a program
\citep{tpp17}. \citet{tpp16} learn a policy with memory, an automaton over histories, because
a rule of the kind studied here cannot say that a push is allowed only after the tests have
run. All of them take the list of rules as given. This paper is about the list.

\section{Limits}\label{sec:limits}

\paragraph{Words are not effects.} A rule matches the words of a command, and the command
does what it does. A forced push can be written with a short option, through an alias, inside
a script, or after a change of directory, and a rule about three words catches none of these.
No analysis of the list can repair that, and this paper does not try to. What it says is that
the list decides what its author thinks it decides, for the patterns that the list has.

\paragraph{The patterns of real harnesses.} We treat words, a wildcard for one word, and a
rest. The rules of Claude Code match the text of a command, and a star in them can stand
inside a word; its rules for files are patterns for paths, in which a part can be matched in
part \citep{claudecodepermissions}. For such patterns usefulness is the inclusion of regular
languages, and the single pass is lost. Rules can also come from several files, of a user, a
project and an organization, that are merged before they are read. The merged list is a list,
and the diagnoses apply to it. Saying which file a dead rule came from, and whose rule killed
it, needs more than this paper has.

\paragraph{No real lists.} The experiments use random lists. We have not collected the lists
that people write, and the shares in Table~\ref{tab:growth} should not be read as rates of
defects in practice. The example of Section~\ref{sec:example} is invented.

\paragraph{Exceptions and accidents.} We report a rule only if it is dead or changes
nothing. A denial that is alive can still lose a part of what it matches to a permission
above it, and whether that is an exception or an accident is not in the list.

\paragraph{Proofs on paper.} The proofs here are short, and none has been checked by a
machine. The script checks their statements on small lists, and exhaustively only by
Lemma~\ref{lem:finite}. That lemma itself is tested only in so far as the comparisons that
rely on it are between an analysis and a definition.

\section{Conclusion}

A permission list is a program, a small one in a small language, and it deserves the warnings
that compilers give for programs. The warnings are cheaper here than there, because commands
are made of words that no list can exhaust. A harness could run the four diagnoses whenever a
rule is added and say at once that the new rule is dead, which rule killed it, and whether
that loosens what the author asked for. It could show the list in normal form, where what is
decided is not hidden behind what was written. Neither needs anything that a compiler does
not already do for a match expression.

\bibliographystyle{ACM-Reference-Format}
\bibliography{refs}

\end{document}
